% retained preamble for sources/alco-235-vertex-models-for-canonical-grothendieck-polynomials-2023
% Assembled by the source-intake line 2026-10-15, per the consumer's six-round error chain spec:
%   mathtools + \DeclarePairedDelimiter, \address, \datereceived/\dateaccepted/\daterevised as no-ops,
%   the abstract environment family, and the ALCO theorem family; author macros kept verbatim.
% Comments are ASCII-only and no orphan metadata arguments remain (delivery rule).
\PassOptionsToPackage{dvipsnames,svgnames,table}{xcolor}
\documentclass[11pt]{article}
\usepackage[a4paper,margin=22mm]{geometry}
\usepackage[no-math]{fontspec}
\setmainfont{DejaVu Serif}
\usepackage{amsmath,amssymb,amsthm,mathtools,graphicx,xcolor,hyperref,xurl,xspace,ifthen,enumerate}
\usepackage[shortlabels]{enumitem}
\setlength{\emergencystretch}{4em}
% ---- publisher-only commands as no-ops (cedram class) ----
\providecommand{\address}[1]{}
\providecommand{\email}[1]{}
\providecommand{\urladdr}[1]{}
\providecommand{\datereceived}[1]{}
\providecommand{\dateaccepted}[1]{}
\providecommand{\daterevised}[1]{}
\providecommand{\datepublished}[1]{}
\providecommand{\subjclass}[1]{}
\providecommand{\keywords}[1]{}
% ---- abstract environment family ----
\providecommand{\abstractname}{Abstract}
\ifdefined\abstract\else\newenvironment{abstract}{\begin{quote}\small\noindent\textbf{\abstractname.}\ }{}\fi
% ---- ALCO theorem family (shared counter, publisher numbering) ----
\newtheorem{theo}{Theorem}[section]
\newtheorem{lemm}[theo]{Lemma}
\newtheorem{coro}[theo]{Corollary}
\newtheorem{prop}[theo]{Proposition}
\newtheorem{defi}[theo]{Definition}
\newtheorem{rema}[theo]{Remark}
\newtheorem{exam}[theo]{Example}
\newtheorem{ques}[theo]{Question}
\newtheorem{conj}[theo]{Conjecture}
% ---- author macros (verbatim) ----
%\usepackage{fourier}
\usepackage[T1]{fontenc}
%\usepackage[a4paper,left=2.25cm,right=2.25cm,top=3cm,bottom=3cm,headsep=1cm]{geometry}
%\setlength{\parskip}{1pt}
\usepackage{hyperref}
\usepackage{cleveref}
%\usepackage{mathtools} %for automatic size of the brackets
%tikz stuff
\usepackage{tikz}\usetikzlibrary{graphs,quotes,fit,positioning,matrix,calc,decorations.markings,angles,decorations.pathmorphing,decorations.pathreplacing}%,matrix,fit,shapes.geometric}
% useful macros
\tikzset{math mode/.style = {execute at begin node=$, execute at end node=$}}%makes all nodes math mode
\tikzset{arrow/.style={postaction={decorate,thick,decoration={markings,mark = at position #1 with {\arrow{>}}}}},arrow/.default=0.5}
\tikzset{invarrow/.style={postaction={decorate,thick,decoration={markings,mark = at position #1 with {\arrow{<}}}}},invarrow/.default=0.5}
\newcommand\bbull[3]{\filldraw[fill=black!35!blue, draw=black] (#1,#2) circle (#3cm);}
\newcommand\ebull[3]{\filldraw[fill=white,draw=black] (#1,#2) circle (#3cm);}
\tikzset{fermionic/.style={thin}}
\tikzset{bosonic/.style={ultra thick}}
\tikzset{gcol/.style={Bo}}
\tikzset{gccol/.style={red}}
\tikzset{Gcol/.style={blue}}
% 
\DeclarePairedDelimiter\autobracket{(}{)}
\newcommand{\br}[1]{\autobracket*{#1}}
%
\definecolor{Bo}{rgb}{1.0, 0.49, 0.0}
\definecolor{Bs}{rgb}{1.0, 0.44, 0.37}
\definecolor{aj}{rgb}{0.8, 0.8, 1.0}
%
\newtheorem{case}{Case}
\newcommand{\sslash}{\mathbin{/\mkern-6mu/}}
\newcommand{\bra}[1]{\left\langle #1\right|}
\newcommand{\ket}[1]{\left|#1\right\rangle}
\newcommand{\braket}[2]{\left< #1 \vphantom{#2}\right|
\! \left. #2 \vphantom{#1} \right>}%
%\usepackage[toc,page]{appendix}
%\usepackage[inline]{showlabels}
\setlength{\multlinegap}{0pt}
\allowdisplaybreaks[1]
%
\renewcommand\ss{\scriptstyle}
\newcommand\sss{\scriptscriptstyle}
\newcommand\ds{\displaystyle}
\def\phid{\phi^\dagger}
\def\psid{\psi^\dagger}
\def\fs{\footnotesize}
\newcommand\g{{g^{(\alpha,\beta)}_{\lambda}}}
\newcommand\Gg{{G^{(\alpha,\beta)}_{\lambda}}}
\newcommand\lb{{\left(}}
\newcommand\rb{{\right)}}
\newcommand\ZZ{{\mathbb Z}}
\newcommand\QQ{{\mathbb Q}}
\newcommand\CC{{\mathbb C}}
\newcommand\PP{{\mathbb P}}
\newcommand\R{{\mathcal {R}}}
\newcommand\LL{{\mathcal {L}}}
\newcommand\jl{{\mathfrak l}}
\newcommand\jt{{\mathfrak t}}


\long\def\junk#1{}







% The abstract text of the source is dropped here: it belongs to the document body, not to the preamble, and the assembler supplies the front matter itself.


% Added for this card: the retained statement uses the thm environment, which the delivered preamble does not declare, so it is aliased to the shared theorem counter.
\newtheorem{thm}[theo]{Theorem}
